\usepackage[no-math]{fontspec}
\setmainfont{DejaVu Serif}
\usepackage{mathtools,proof,tikz,csquotes,needspace,enumerate,xspace,array,xcolor}
\newcommand{\faBible}{\raisebox{-0.1254em}{\includegraphics[height=1em]{author-bible-glyph.pdf}}}
\usepackage{hyperref}
\usepackage[capitalise]{cleveref}
\usetikzlibrary{calc,positioning}
\definecolor{lipicsYellow}{RGB}{255,204,0}
\definecolor{lipicsGray}{gray}{0.4}
\tikzstyle{m}=[fill=white,inner sep=1,scale=0.75]
\tikzstyle{n}=[fill=lipicsYellow!50,inner sep=1,rounded corners]
\tikzstyle{b}=[lipicsGray,thick]
\newtheorem{theorem}{Theorem}
\newtheorem{lemma}[theorem]{Lemma}
\newtheorem{proposition}[theorem]{Proposition}
\newtheorem{corollary}[theorem]{Corollary}
\theoremstyle{definition}
\newtheorem{definition}[theorem]{Definition}
\newtheorem{example}[theorem]{Example}
\newtheorem{remark}[theorem]{Remark}
\newtheorem{notation}[theorem]{Notation}
\newtheorem{problem}[theorem]{Problem}
\providecommand{\todo}[2][]{}
% BEGIN unchanged author macros.tex
% !TEX root = CerdaMS.tex

% ========== Comments macros ==========

\newcommand{\comas}	[1]	{\todo[color=cyan]			{Alex: #1}}
\newcommand{\comasi}[1]	{\todo[color=cyan,inline]	{Alex: #1}}
\newcommand{\comrc}	[1]	{\todo[color=orange]		{Rémy: #1}}
\newcommand{\comrci}[1]	{\todo[color=orange,inline]	{Rémy: #1}}
\newcommand{\comgm}	[1]	{\todo[color=green]			{Giulio: #1}}
\newcommand{\comgmi}[1]	{\todo[color=green,inline]	{Giulio: #1}}


% ========== Text macros ==========

\newcommand{\bsub}{\begin{enumerate}[(i)]}
\newcommand{\esub}{\end{enumerate}}

\NewDocumentCommand \lam		{}		
	{\texorpdfstring{\ensuremath{\lambda}}{lambda}}
\NewDocumentCommand \lamI		{}		
	{\texorpdfstring{\ensuremath{\lambda\comb{I}}}{lambda-I}}
\NewDocumentCommand \lamK		{}		
	{\texorpdfstring{\ensuremath{\lambda\comb{K}}}{lambda-K}}

% Emphasizing defined terms
\NewDocumentCommand \defemph 	{m} 	{\emph{#1}}


% ========== Basic stuff ==========

% Defining objects

\NewDocumentCommand \eqdef		{}		{ \coloneqq }
\NewDocumentCommand \eqbnf		{}		{ \mathrel{\Coloneqq} }
\NewDocumentCommand \eqcobnf	{}		
	{ \mathrel{\Coloneqq}_{\textrm{co-ind.}} }
\NewDocumentCommand \bnf		{ommo}
	{\[	\IfValueT{#1}{#1 \quad \ni \quad}
		#2 \eqbnf #3
		\IfValueT{#4}{\quad (#4)}
	\]}
\NewDocumentCommand \cobnf		{ommo}
	{\[	\IfValueT{#1}{#1 \quad \ni \quad}
		#2 \eqcobnf #3
		\IfValueT{#4}{\quad (#4)}
	\]}

% Set theory

\NewDocumentCommand \set 		{smo}	
	{\IfBooleanTF{#1}
		{ \left\{ #2 \IfValueT{#3}{ \ \middle|\ #3 } \right\} }
		{ \{ #2 \IfValueT{#3}{ \ |\ #3 } \} }
	}

\NewDocumentCommand \subf		{}		{ \subseteq_{\mathrm{f}} }
\RenewDocumentCommand \emptyset {}		{ \varnothing }
\RenewDocumentCommand \setminus {}		{ - }
\NewDocumentCommand	\parts		{m}		{ \cramped{2^{#1}} }
\NewDocumentCommand \bigdunion	{}		{ \coprod }

% Predefined sets

\NewDocumentCommand \Nat 		{}		{ \mathbb{N} }
\NewDocumentCommand \Var		{}		{ \mathcal{V} }

% Rewriting

\NewDocumentCommand \red 			{oo}	
	{ \rightarrow \IfValueTF{#1}
		{ \IfValueTF{#2}{ _{#1}^{#2} }{ _{#1} } }
		{ \IfValueT{#2}{ ^{#2} } } 
	}
\NewDocumentCommand \redd 			{oo}	
	{ \mathrel{\rightarrow\mathrel{\mkern-14mu}\rightarrow} \IfValueTF{#1}
		{ \IfValueTF{#2}{ _{#1}^{#2} }{ _{#1} } }
		{ \IfValueT{#2}{ ^{#2} } } 
	}


% ========== λ(I)-calculus ==========


% Sets of λ-terms
\NewDocumentCommand \Lam		{}		{ \Lambda }
\NewDocumentCommand \LamI		{}		{ \Lambda_\comb{I} }

\NewDocumentCommand \subst 		{om}	
	{[ #2 / \IfValueTF{#1}{#1}{x}] }

% Combinators
\NewDocumentCommand \comb		{m}		{ \mathsf{#1} }
\NewDocumentCommand \Om			{}		{ \comb{\Omega} }
\NewDocumentCommand \Bible 		{o} 		
	{ \text{\scalebox{0.8}{\faBible}}_{\,\IfValueTF{#1}{#1}{l}} }

% \ensuremath{\beta}-reduction
\NewDocumentCommand \bred 		{} 		{ \red[\ensuremath{\beta}] }
\NewDocumentCommand \breds 		{} 		{ \redd[\ensuremath{\beta}] }


% ========== Resource calculus ==========

% Sets of resource expressions
\NewDocumentCommand \rterms		{ss}
	{\mathord{
		\IfBooleanTF{#1}{ \IfBooleanTF{#2}{(!)}{!} }{}
		\Delta_{\comb{I}}
	}}

% Sets of finite sums of resource expressions
\NewDocumentCommand \frsums 	{oss}
	{ \IfBooleanTF{#2}
		{ \IfBooleanTF{#3}
			{ \Nat[ \rterms** \IfValueT{#1}{(#1)} ] }
			{ \Nat[ \rterms* \IfValueT{#1}{(#1)} ] } }
		{ \Nat[ \rterms \IfValueT{#1}{(#1)} ] }
	}

% Empty sum
\NewDocumentCommand \esum		{m}		{ \rs{0}_{#1} }

% Formatting resource terms, bags, sums, and sums of bags
\NewDocumentCommand \rt			{m}		{ #1 }
\NewDocumentCommand \rb			{m}		{ \bar{#1} }
\NewDocumentCommand \rs			{m}		{ \mathbf{#1} }
\NewDocumentCommand \rbs		{m}		{ \bar{\mathbf{#1}} }

% Memory substitution
\NewDocumentCommand \msubst 	{mo}	
	{ \left\lbrace #1 \middle/ \IfValueTF{#2}{#2}{x} \right\rbrace }

% Resource substitution
\NewDocumentCommand \rsubst 	{smo}	
	{ \left\langle \IfBooleanTF{#1}{ #2 }{ \rb{#2} } \middle/ 
	\IfValueTF{#3}{ #3 }{ x } \right\rangle }

% Resource reduction
\NewDocumentCommand \resource	{}		{ \mathsf{r} }
\NewDocumentCommand \rred 		{} 		{ \red[\resource] }
\NewDocumentCommand \rredr 		{} 		{ \red[\resource][?] }
\NewDocumentCommand \rreds 		{} 		{ \redd[\resource] }

\NewDocumentCommand \nf 		{}		{ \mathrm{nf} }
\NewDocumentCommand \size		{m}		{ \mathsf{size}(#1) }
\NewDocumentCommand \sumsize	{m}		{ \mathsf{ssize}(#1) }

\NewDocumentCommand \supp		{sm}
	{ | \IfBooleanTF{#1}{#2}{\rs{#2}} | }

% Multiset ordering

\NewDocumentCommand \msetlt		{}		{ \prec_\Nat }
\NewDocumentCommand \msetgt		{}		{ \succ_\Nat }


% ========== Böhm-like trees and approximations ==========

\NewDocumentCommand \memory		{}		{ \mathsf{m} }

\NewDocumentCommand \fv 		{}		{ \mathrm{fv} }
\NewDocumentCommand \pfv 		{}		{ \mathrm{pfv} }

% Trees

\NewDocumentCommand \BT 		{m}		{ \mathrm{BT}(#1) }
\NewDocumentCommand \MT			{om}	
	{ \mathrm{OT}\IfValueT{#1}{^{#1}}(#2) }
\NewDocumentCommand \MTrees		{}		{ \mathcal{OT} }

\NewDocumentCommand \Appset		{}		{ \mathcal{A}_\memory }
\NewDocumentCommand \Appof		{m}		{ \mathrm{App}_\memory(#1) }
\NewDocumentCommand \dirapp		{m}		{ \omega_\memory(#1) }
\NewDocumentCommand \MTle		{}		{ \sqsubseteq }
\NewDocumentCommand \MTlub		{}		{ \sqcup }
\NewDocumentCommand \bigMTlub	{}		{ \bigsqcup }

% The bijection Ohana trees ↔ Sets of approximants
\NewDocumentCommand \mtapp		{m}		{ {\cal A}(#1) }
\NewDocumentCommand \appmt		{m}		{ {\cal T}(#1) }

% Taylor

\NewDocumentCommand \TE 		{ssm}	
	{ \mathcal{T}_\memory^{
		\IfBooleanT{#1}{\IfBooleanTF{#2}{!}{\bullet}}
	} (#3) }
\NewDocumentCommand \TMT		{ssm}
	{ \IfBooleanTF{#1}
		{ \IfBooleanTF{#2}{ \TE**{\MT{#3}} }{ \TE*{\MT{#3}} } }
		{ \TE{\MT{#3}} } 
	}
\NewDocumentCommand \NFT		{ssm}
	{ \nf( \IfBooleanTF{#1}
		{ \IfBooleanTF{#2}{ \TE**{#3} }{ \TE*{#3} } }
		{ \TE{#3} } 
	) }

\NewDocumentCommand \TEin		{}		{ \in^\bullet }
\NewDocumentCommand \TElt		{}		{ \subseteq^\bullet }
\NewDocumentCommand \TElub		{}		{ \cup^\bullet }
\NewDocumentCommand \bigTElub	{}			
	{ \mathop{\bigcup\nolimits^{\mathrlap{\bullet}}}\displaylimits }

\NewDocumentCommand \resourceof	{sm}	
	{ [\![#2]\!]_\resource^{\IfBooleanT{#1}{!}} }

% Associated theories

\NewDocumentCommand \cB			{}		{ \mathcal{B} }
\NewDocumentCommand \cT			{}		{ \mathcal{T} }
\NewDocumentCommand \cM			{}		{ \mathcal{O} }
\NewDocumentCommand \cX			{}		{ \mathcal{X} }
\NewDocumentCommand \cY			{}		{ \mathcal{Y} }
\NewDocumentCommand \model		{m}		{ \mathcal{#1} } 
% find a font for models


% END unchanged author macros.tex

\providecommand{\cal}{\mathcal}
